{\footnotesize
\setlength{\tabcolsep}{4pt}
\renewcommand{\arraystretch}{1.12}
\begin{longtable}{@{}>{\raggedright\arraybackslash}p{0.33\linewidth}>{\raggedright\arraybackslash}p{0.64\linewidth}@{}}
\caption{Archive-by-archive comparison at the pinned ProveIt commit.}\\
\toprule Archive & Distinct contribution and validation boundary\\\midrule
\endfirsthead
\toprule Archive & Distinct contribution and validation boundary\\\midrule
\endhead
\file{Below_U208_Research_2026-10-29.zip} & Comparison and bounded local, source, and offset searches in the U208 family. No smaller complete machine is supplied. \\ \addlinespace[5pt]
\file{ConPA_below208_research_2026-10-29.zip} & Source permutations, two-row screens, and a logical obstruction in the U208 family. No sub-208 construction is established. \\ \addlinespace[5pt]
\file{PA_consistency_TM_research_2026-11-13.zip} & Fourteen-table census. A 3,573-site source has 157 grammar nodes, but no new 196- or 195-state physical table is supplied; a residual-MP proposal also remains a source-level experiment. \\ \addlinespace[5pt]
\file{U208_ConPA_2026-10-29.zip} & Original U208 class A, using the 172/211 entry-history fold. Other presentations of the same table are deduplicated. \\ \addlinespace[5pt]
\file{U208_ConPA_Conditional_Context_2026-10-29.zip} & Original U208 classes A and B, presented through a conditional-history invariant. The proof presentation does not create an additional graph. \\ \addlinespace[5pt]
\file{U208_ConPA_Entry_Context_2026-10-29.zip} & Original U208 classes A and B, with an entry-context proof presentation. \\ \addlinespace[5pt]
\file{U208_ConPA_History_2026-10-29.zip} & Original U208 classes A and B, with a predecessor-history proof presentation. \\ \addlinespace[5pt]
\file{conpa_tm_comparison_and_reduction_search.zip} & Two experimental 196-state MP variants with bounded frame checks. A rejected Shift edit halts after 43,563 blank steps; these proof scopes are distinguished from the later uniform Combined check. \\ \addlinespace[5pt]
\file{conpa_tm_comparison_and_search.zip} & An A/D register-swap 196-state table. A rotated 195-state proposal fails validation; an unmodified scanner shortcut halts after 699 blank steps. \\ \addlinespace[5pt]
\file{conpa_tm_study_2026-11-13.zip} & Combined196: H-only Shift and MP result in H, 3,788 source sites, 4,928 uniform cases, and 57,922-step startup. A direct-negation variant has 197 states; a repaired 157-node source needs 207 physical states. \\ \addlinespace[5pt]
\file{pa196_research_2026-11-13.zip} & A 157-node source without a corresponding accepted binary table, a lost-helper obstruction, and negative bounded two-row synthesis results. \\ \addlinespace[5pt]
\file{pa208_research_2026-10-29.zip} & Eighteen additional pairwise nonisomorphic 208-state helper layouts. They extend the two original U208 graph classes to twenty designated classes. \\ \addlinespace[5pt]
\file{pa_below196_research.zip} & An earlier 23-archive census with fifteen selected 196-state graphs. Equivalent gate spellings, a restricted implication-matrix minimum, and bounded physical-fold searches do not give a sub-196 machine. \\ \addlinespace[5pt]
\file{pa_consistency_comparison_2026_11_13.zip} & Quantifier-head splitting with Hilbert certificates and twelve combined source/table experiments. The best combined negation/MP layout has 197 states; further gate and padding searches do not improve it. \\ \addlinespace[5pt]
\file{pa_consistency_comparison_and_new196_2026-11-13.zip} & Three diagonal-calculus 196-state layouts: diagonal, separate head, and scanner-oriented separate head. A hybrid passes frame contracts but its tested starts escape outward. A short-scanner 195-state proposal fails its clean seed. \\ \addlinespace[5pt]
\file{pa_consistency_reduction.zip} & Rebracketed196 and H-only Shift196, plus integrated witness-compiler experiments costing 203 or 206 states. \\ \addlinespace[5pt]
\file{pa_consistency_research.zip} & A 158-node source compiles to 196 states after an eight-zero scanner repair. The shorter uncorrected scanner is not an accepted reduction. \\ \addlinespace[5pt]
\file{pa_consistency_research_2026-10-29.zip} & A source reduction from 3,909 to 3,904 sites increases the relevant DAG allocation from 197 to 200 nodes. No smaller literal table is obtained. \\ \addlinespace[5pt]
\file{pa_consistency_research_2026-11-13 (2).zip} & A further distinct 196-state twelve-rotation layout; a 198-state diagonal-vacuity layout; a halt-at-successful-MP command theorem; and reduced joint register abstractions. \\ \addlinespace[5pt]
\file{pa_consistency_research_2026-11-13.zip} & A 199-state diagonal-vacuity implementation and a 197-state negation peephole. Neither satisfies the strict target. \\ \addlinespace[5pt]
\file{pa_tm_comparison_and_reductions.zip} & Two rotated 196-state layouts and an equivalent-gate compiler with 3,773 source sites and 198 physical states. \\ \addlinespace[5pt]
\file{pa_tm_extended_research_2026-11-13.zip} & Four 196-state closing-decrement variants, with their stated constrained symbolic and startup checks. \\ \addlinespace[5pt]
\file{pa_tm_research_2026-11-13 (2).zip} & Zerojump196 uses the source-site-26 zero-register invariant: 4,959 uniform cases and 58,402-step startup. A 195-state fold fails on the actual blank run. Base-8 and disequality-parity proposals have incomplete implementation or startup obligations. \\ \addlinespace[5pt]
\file{pa_tm_research_2026-11-13.zip} & An alternative 196-state layout with maximum source-path length 36, full symbolic instruction checks, and exact startup validation. \\ \addlinespace[5pt]
\file{pa_unary_power_cover.zip} & A unary-cover theory and proof-stream proposal. No complete binary compilation or sub-196 transition table is supplied. \\ \addlinespace[5pt]
\file{u208_comparison_research_2026-10-29.zip} & A 207-state U208-family proposal rejected by a concrete primitive-frame counterexample. \\ \addlinespace[5pt]
\file{u208_research_2026-10-29.zip} & Original U208 class B as the primary table and class A as an alternative; these repeat graph classes in the companion presentations. \\ \addlinespace[5pt]
\bottomrule
\end{longtable}
}
